\section{Uniform estimates and the Kac approximation}\label{sec:kac}

The next step relates a physical long-range attraction to the transform
\(\T_\sigma\) in \eqref{eq:transform}.  Throughout this section, \(t\) satisfies
\eqref{eq:t-choice} and is fixed.  All uniform estimates concern the physical
class \eqref{eq:physical-class} at this fixed \(t\); they need not remain
uniform as \(t\) tends to infinity.  The approximation is a form of the
Lebowitz--Penrose limit, whose extension to superstable reference interactions
is classical \cite{LebowitzPenrose1966,GerardiEtAl1973}.  We give the proof
because we will need constants independent of the attraction ranges, including
in auxiliary systems with some positive interactions removed.

\subsection{Local moments and continuity in the attractive tail}

Recall the cells \(C_i\) of side \(b\), their occupation numbers \(n_i\), and
the attractive interaction bounds \(J_{ij}\) from \eqref{eq:cell-row}.
In a finite cube, consider the grand canonical measure with activity \(e^{th}\).
We also allow two modifications: positive core interactions may be deleted
between boxes that are unions of whole cells, and negative pair interactions
may be multiplied by measurable symmetric factors in \([0,1]\).
Every original cell retains its internal core interaction.  Consequently,
the lower bound \eqref{eq:superstable} continues to hold.

\begin{proposition}[Uniform local moments]\label{prop:moments}
Fix a compact set \(\mathcal K\subset\Omega\).
For all potentials in \eqref{eq:physical-class}, all finite cubes, and all
the modified systems just described, there is a constant \(C_{\mathcal K}\)
such that
\[
 \mathbb E_{\beta,h}[n_i n_j]\le C_{\mathcal K}
 \qquad ((\beta,h)\in\mathcal K)
\]
for every pair of cells meeting the cube.  The constant is independent of
the number and ranges of the attractive summands.
\end{proposition}

\begin{proof}
Condition on the configuration \(Y\) outside \(C_i\), and put
\[
 A_i(Y)=\sum_{j\ne i}J_{ij}n_j,\qquad c=\frac{\beta A_0}{4}.
\]
The interactions involving the \(n\) particles inside \(C_i\) are bounded
below by
\[
 (A_0-J_{ii}/2)n^2-B_0n-nA_i.
\]
The empty cell contributes one to the conditional partition function.
Its volume is at most \(b^3\), also when cropped by the cube boundary.
It follows that
\begin{align*}
 \mathbb E[e^{c n_i^2}\mid Y]
 &\le
 \sum_{n\ge0}\frac{b^{3n}}{n!}
 \exp\!\left(cn^2+thn-\beta(A_0-J_{ii}/2)n^2
                  +\beta B_0n+\beta nA_i\right)\\
 &\le C_1\exp(\beta A_i^2/A_0).
\end{align*}
Indeed,
\(\beta nA_i\le\beta A_0n^2/4+\beta A_i^2/A_0\);
after this estimate the remaining negative quadratic coefficient is at least
\[
 \beta(A_0/2-M/2)\ge 7\beta A_0/16.
\]
The resulting one-variable series is bounded uniformly on \(\mathcal K\).
The conditional bound still depends on \(Y\).  We remove that dependence
by a finite-volume bootstrap.

Let \(F_i=\mathbb E[e^{c n_i^2}]\).  These numbers are finite before the
bootstrap: multiplying the grand canonical integrand by \(e^{cn_i^2}\)
leaves a strictly negative quadratic in all cell counts, by
\eqref{eq:superstable}.  Cauchy--Schwarz with weights \(J_{ij}\) gives
\[
 A_i^2\le M\sum_{j\ne i}J_{ij}n_j^2.
\]
If \(M>0\), set \(\theta_{ij}=J_{ij}/M\) and
\(\alpha=4(M/A_0)^2<1\).  Then $c\alpha=\beta M^2/A_0$. The row weights sum to at most one;
place any missing weight on a value zero. Convexity of the exponential and then
concavity of \(x^\alpha\) give
\[
 F_i\le C_1\left(1+\sup_j F_j^\alpha\right).
\]
There are finitely many cells, so their maximum is bounded by a constant
depending only on \(C_1,\alpha\).  When \(M=0\), the conditional estimate
already gives \(F_i\le C_1\).  Since \(\beta\) is bounded away from zero on
\(\mathcal K\), this proves uniform second moments.  A final
Cauchy--Schwarz inequality proves the assertion.
\end{proof}

This is a direct version of the local-occupancy estimates associated with
superstability \cite{Ruelle1970}.  Its extra uniformity will control both
the infinite tail and the boundary errors in the Kac comparison.

\begin{proposition}[Attractive perturbations]\label{prop:tail}
Fix a compact set \(\mathcal K\subset\Omega\).
For two of the preceding finite-volume systems differing by a negative
interaction, let \(J^{\mathrm{diff}}_{ij}\) be cellwise suprema of its
absolute value.  If the linear interpolation stays within the same
attractive mass bound, then
\begin{equation}\label{eq:pressure-interpolation}
 |\log\Xi(1)-\log\Xi(0)|
 \le \frac{\beta C_{\mathcal K}}2\sum_{i,j}J^{\mathrm{diff}}_{ij}
 \qquad ((\beta,h)\in\mathcal K).
\end{equation}
In particular, for
\[
 \varphi=\psi_t-t\sum_{j\ge1}\sigma_jw_{R_j},
 \qquad
 \varphi_n=\psi_t-t\sum_{j\le n}\sigma_jw_{R_j},
\]
there is a range-independent constant \(C'_{\mathcal K}\) such that
\begin{equation}\label{eq:tail-continuity}
 \|Q_\varphi-Q_{\varphi_n}\|_{\mathcal K}
 \le C'_{\mathcal K}\sum_{j>n}\sigma_j .
\end{equation}
\end{proposition}

\begin{proof}
Differentiate the logarithm of the grand partition function along the
interpolation.  The absolute derivative is at most
\[
 \frac{\beta}{2}\sum_{i,j}J^{\mathrm{diff}}_{ij}\,
                 \mathbb E[n_i n_j].
\]
Including particle self pairs only enlarges this bound.  Differentiation
under the integrals and sums is valid by the common negative quadratic
in the finitely many cell counts.  Proposition~\ref{prop:moments} and
integration over the interpolation parameter prove
\eqref{eq:pressure-interpolation}.

For the tail difference, each row sum is at most
\(c_b t\sum_{j>n}\sigma_j\), and the number of cells is at most
\((L/b+2)^3\).  Divide \eqref{eq:pressure-interpolation} by \(tL^3\),
use the pressure limit in Proposition~\ref{prop:thermodynamics}, and let
\(L\to\infty\).  The bounds are uniform on \(\mathcal K\), proving
\eqref{eq:tail-continuity}.
\end{proof}

\subsection{One long-range attraction}

\begin{proposition}[Physical Kac approximation]\label{prop:kac}
Let
\[
 \varphi=\psi_t-t\sum_{j=1}^k\sigma_jw_{R_j}
\]
have finitely many attractive summands, with the empty sum allowed.
Fix \(\sigma\ge0\) such that \(\sum_{j=1}^k\sigma_j+\sigma\le A+1\).
Then
\begin{equation}\label{eq:kac-limit}
 Q_{\varphi-t\sigma w_R}\longrightarrow\T_\sigma Q_\varphi
 \quad\text{as }R\to\infty
\end{equation}
uniformly on compact subsets of \(\Omega\).
\end{proposition}

\begin{proof}
Fix the reference potential, and let \(R_0\) bound its interaction range.
All constants depending on \(R_0\) are fixed during this proof.  Partition
a cube of side \(L=mu\) into \(m^3\) boxes of side
\[
 u=b\left\lfloor R^{1/2}/b\right\rfloor.
\]
Translate the \(b\)-cell lattice, if necessary, to align it with the
outside cube; all preceding estimates are invariant under this
translation.  The \(u\)-boxes are then unions of whole \(b\)-cells.
For each \(R\), we
first let \(m\to\infty\); only afterward do we let \(R\to\infty\).
In particular \(u\to\infty\), \(u/R\to0\), and eventually \(u\gg R_0\).

Write \(B_z\) for the box centered at \(z\), and define
\[
 K^+_{zz'}=\sup_{x\in B_z,\ y\in B_{z'}}w_R(x-y),
 \qquad
 K^-_{zz'}=\bigl(w_R(z-z')-Cu/R^4\bigr)_+.
\]
Choose the fixed constant \(C\) large enough that \(K^-_{zz'}\) is at most
the corresponding infimum.  The \(C^1\) bound on \(w\) and its compact
support imply, for some \(\eta_R=O(u/R)\to0\),
\begin{equation}\label{eq:kac-rows}
 u^3\sum_{z'}K^+_{zz'}\le1+\eta_R.
\end{equation}
If \(z\) is farther than \(R+3u\) from the outside cube boundary, then
\begin{equation}\label{eq:kac-lower-row}
 \sum_{z'}K^-_{zz'}\ge(1-\eta_R)u^{-3}.
\end{equation}
Here the sums are over the finite center grid in the cube.  To check these
estimates, compare \(u^3\sum_{z'}w_R(z-z')\) with \(\int w_R=1\).
The oscillation in each grid cube is \(O(uR^{-4})\), and only
\(O((R/u)^3)\) cubes meet the support enlarged by \(3u\).  The total
error after multiplication by \(u^3\) is \(O(u/R)\).
The same argument controls the suprema and the truncated lower kernel.
Also \(K^\pm_{zz'}\le C R^{-3}\).

\paragraph{Upper bound.}
Delete all positive core interactions between the \(u\)-boxes.  This
increases the partition function.  Remove next the reference attractive
interactions between these boxes, keeping the new Kac attraction in
place.  By Proposition~\ref{prop:tail}, this second change costs at most
\(C_{\mathcal K,R_0}L^3/u\) in the logarithm.  Indeed only original cells
within \(R_0+2b\) of a box boundary have a nonzero difference row, and
there are \(O_{R_0}(L^3/u)\) such cells.  Proposition~\ref{prop:moments}
applies throughout this interpolation.

Write \(n_z\) for the count in the box centered at \(z\).
The long-range pair sum is at most
\[
 \frac12\sum_{z,z'}K^+_{zz'}n_zn_{z'}
 \le \frac{1+\eta_R}{2u^3}\sum_z n_z^2,
\]
where the last step uses symmetry and
\(n_zn_{z'}\le(n_z^2+n_{z'}^2)/2\).
The first bound includes self pairs, which is harmless for this upper
estimate of an attractive Boltzmann factor.  The remaining reference
interactions are confined to individual boxes.  Summing the assignments
of particle labels cancels the global canonical factorial and gives
\begin{equation}\label{eq:kac-upper}
 \Xi_{mu}^{\varphi-t\sigma w_R}
 \le
 \exp(C_{\mathcal K,R_0}L^3/u)
 \left[
 \sum_{n\ge0}e^{thn}Z_{u,n}^{\varphi}(\beta)
 \exp\!\left(\frac{\beta t\sigma(1+\eta_R)}{2u^3}n^2\right)
 \right]^{m^3}.
\end{equation}

The finite-volume bound \eqref{eq:finite-coercivity} and
\(\kappa>10t(A+1)\) leave the effective energy coefficient
\[
 \frac{\kappa}{(1+2b/u)^3}-\frac{t\sigma(1+\eta_R)}2>0
\]
uniformly away from zero for all sufficiently large $R$. Consequently
counts \(n/u^3>D\), for sufficiently large \(D\), are suppressed by a
uniform negative quadratic.  The sum of
these high-count terms is exponentially small compared with the empty
term.  Among the remaining \(O(u^3)\) terms, every maximizing subsequence
has a further subsequence with \(n/u^3\to d\in[0,D]\).
The all-sequence limit \eqref{eq:canonical-limit} then gives
\[
 \lim_{R\to\infty}\frac1{tu^3}
 \log\!\left[
 \sum_{n\ge0}e^{thn}Z_{u,n}^{\varphi}(\beta)
 \exp\!\left(\frac{\beta t\sigma(1+\eta_R)}{2u^3}n^2\right)
 \right]
 =\sup_{d\ge0}\left\{hd-H_\varphi(\beta,d)
                         +\frac{\beta\sigma d^2}{2}\right\}.
\]
The logarithm of the number of low-count terms divided by \(u^3\)
vanishes.  This proves the equality by the preceding upper subsequence
argument and by taking \(n=\lfloor du^3\rfloor\) for the lower bound.
Taking first \(m\to\infty\) in \eqref{eq:kac-upper} and then \(R\to\infty\)
proves the desired upper estimate.

\paragraph{Lower bound.}
Place a cube of side \(s=u-2\) at each of the same centers, and require
exactly \(n=\lfloor ds^3\rfloor\) particles in every smaller cube,
where \(d\ge0\) is fixed.  The gaps have width two.  Thus all reference
interactions between distinct smaller cubes are nonpositive and can be
discarded in a lower bound on the partition function.
Let \(m_{\mathrm{int}}\) be the number of centers satisfying the
interior condition in \eqref{eq:kac-lower-row}.  The Kac pair sum obeys
\[
 \sum_{\ell<j}w_R(x_\ell-x_j)
 \ge \frac12m_{\mathrm{int}}n^2u^{-3}(1-\eta_R)
                      -\frac12m^3nCR^{-3}.
\]
The last term removes the particle self pairs included in the diagonal
box contribution.  The count constraint and the canonical factorials
therefore give
\begin{align*}
 \Xi_{mu}^{\varphi-t\sigma w_R}
 \ge{}&
 \left(e^{thn}Z_{s,n}^{\varphi}(\beta)\right)^{m^3}\\
 &{}\times
 \exp\!\left\{\frac{\beta t\sigma}{2}
 \left[m_{\mathrm{int}}n^2u^{-3}(1-\eta_R)-m^3nCR^{-3}\right]\right\}.
\end{align*}
For fixed \(R\), \(m_{\mathrm{int}}/m^3\to1\) as \(m\to\infty\).
Next \(R\to\infty\) gives \(s/u\to1\), \(n/s^3\to d\), and a vanishing
self-pair error.  By \eqref{eq:canonical-limit}, the resulting pressure
lower bound is
\[
 hd-H_\varphi(\beta,d)+\frac{\beta\sigma d^2}{2}.
\]
Take the supremum over \(d\ge0\).

We have proved pointwise convergence in \eqref{eq:kac-limit}.
The functions on both sides are finite and convex on the open convex
set \(\Omega\), by Propositions~\ref{prop:thermodynamics}
and~\ref{prop:transform}.  Pointwise convergence of finite convex
functions to a finite convex function is uniform on compact subsets,
which completes the proof.
\end{proof}

The proposition uses boxes only in its upper and lower estimates.
Both sides of \eqref{eq:kac-limit} are built from the full radial
reference interaction.  At every later stage we will fix the finite
reference first and then choose a sufficiently large new range; no
uniform rate over changing reference ranges is required.

\section{An initial triple of supporting slopes}\label{sec:triple}

We now choose \(t\) so that the ideal transform \(\T_AQ_{\psi_t}\)
has three separated groups of supporting slopes. The core permits
zero-energy clusters of sizes \(m=1,2,3,4\). After one root is positioned,
each of the other \(m-1\) particles has a relative position in a region
of volume proportional to \(e^{-3t}\). With activity \(e^{th}\), the
leading exponent per cluster, divided by \(t\), is therefore
\[
 mh-3(m-1).
\]
This expression is affine in \(m\), so its maximum over the four allowed
sizes occurs at \(m=1\) or \(m=4\). Using the limiting packing density
\(p\), together with the empty configuration, gives the three branches
in the next proposition. The proof below controls configurations that
pay a positive core energy as well.

\subsection{The free-cluster limit}

Write \(q_t=Q_{\psi_t}\).  In this section we temporarily let
\(t\to\infty\); one value will be fixed at its end.
The packing density \(p\) and \(A=3/(2p)\) retain their definitions
from Section~\ref{sec:model}.

\begin{proposition}[Reference pressure]\label{prop:free-pressure}
Uniformly on compact subsets of \(\Omega\),
\begin{equation}\label{eq:free-pressure-limit}
 q_t(\beta,h)\longrightarrow
 q_*(h):=\max(0,ph,4ph-9p).
\end{equation}
\end{proposition}

\begin{proof}
Call a configuration free if every pair distance below one lies in
\([e^{-t},(1+\delta)e^{-t}]\).  Join points at distance below one.
A two-edge path has endpoint distance at most
\(2(1+\delta)e^{-t}<1\); freeness makes the endpoints adjacent.
Thus connected components are complete graphs.
The geometric estimate in Lemma~\ref{lem:shell-graph} bounds
their sizes by four.  Representatives from different components
are separated by at least one, so their number is at most \(M_L\).

Let \(V=L^3\), and let \(\Xi_{\mathrm{free}}\) denote the grand
configurational integral restricted to free configurations.  Root
each component and integrate its root over volume \(V\).
Every other point in the component lies in a ball of volume at most
\(c_0e^{-3t}\) about the root.  The labeled block enumeration,
followed by division by the canonical factorial, gives the upper bound
\begin{equation}\label{eq:free-enumeration}
 \Xi_{\mathrm{free}}\le
 \sum_{k=0}^{M_L}\frac{(V\Theta_t)^k}{k!},
 \qquad
 \Theta_t=\sum_{m=1}^4
 \frac{e^{thm}(c_0e^{-3t})^{m-1}}{(m-1)!}.
\end{equation}
For completeness, ordered component sizes give a factor
\(1/(k!\prod m_i!)\); allowing any of the \(m_i\) labels to be a
root replaces \(m_i!\) by \((m_i-1)!\) and only increases the count.
Ignoring separation between the roots also enlarges the integral.

A bad edge is a pair at distance below one but outside the allowed
shell.  If there are \(B\) bad edges, their endpoint set \(S\) has
size at most \(2B\), its complement is free, and the core energy
is at least \(t^2B\).  Hence, pointwise on labeled configurations,
\[
 e^{-\beta U^{\psi_t}}
 \le \sum_{S\subset\{1,\ldots,N\}}
        e^{-\beta t^2|S|/2}\,
        {\bf1}_{\{\text{the complement of }S\text{ is free}\}}.
\]
The term corresponding to the actual bad-edge endpoint set proves
this inequality; zero Boltzmann factors cause no difficulty.
Summing over particle numbers and integrating the removed particles
gives
\begin{equation}\label{eq:bad-particle-factor}
 \Xi_L^{\psi_t}
 \le \exp\!\left(Ve^{th-\beta t^2/2}\right)\Xi_{\mathrm{free}}.
\end{equation}

Now
\[
 \frac1t\log\Theta_t
 \longrightarrow
 \max_{1\le m\le4}\{mh-3(m-1)\}
 =\max(h,4h-9).
\]
The convergence is uniform for bounded \(h\).  For \(u=k/V\),
the factorial estimate \(k!\ge(k/e)^k\) bounds the remaining
normalized logarithm in \eqref{eq:free-enumeration} by
\(u(1-\log u)/t\), with value zero at \(u=0\).
It is uniformly \(o(1)\) for bounded \(u\) as \(t\to\infty\).
There are only \(M_L+1=O(V)\) terms.
Take \(L\to\infty\) first and use \(M_L/V\to p\).
The extra pressure in \eqref{eq:bad-particle-factor} is
\(e^{th-\beta t^2/2}/t\), which tends to zero uniformly on compact
subsets of \(\Omega\).  This proves the upper limit in
\eqref{eq:free-pressure-limit}.

For the lower limit, fix \(\ell>0\) and \(\eta>0\).  Dilate an optimal
\(M_\ell\)-packing by \(1+\eta\), and repeat it with cubic period
\[
 P=(1+\eta)(\ell+2).
\]
In a cube of side \(mP\), the \(k=m^3M_\ell\) resulting centers have
mutual distances at least \(1+\eta\) and a fixed positive margin
from the outside boundary.  At each center put either one particle
or four particles, using the same choice \(m_0\in\{1,4\}\) everywhere.
The root ranges over a ball of fixed volume \(v_0>0\), small compared
with the separation margin.  For \(m_0=4\), put the other particles
in relative balls of volume \(v_1e^{-3t}\) about the three other
vertices of a regular tetrahedron of side
\((1+\delta/2)e^{-t}\).  A sufficiently small fixed \(v_1>0\)
puts every internal distance in the shell.  For all large \(t\),
distances between clusters exceed one.

All these configurations have zero core energy.  Assign unordered
sets of \(m_0\) labels to the \(k\) centers and choose one deterministic
ordering of each set for the root and vertices.  Different assignments
give disjoint integration regions, and the relative coordinates have
Jacobian one.  Thus
\[
 Z_{mP,m_0k}^{\psi_t}(\beta)
 \ge
 \left(\frac{v_0(v_1e^{-3t})^{m_0-1}}{m_0!}\right)^k .
\]
First let \(m\to\infty\), then \(t\to\infty\).
The contribution to \(q_t\) has lower limit
\[
 \frac{M_\ell}{P^3}\bigl(m_0h-3(m_0-1)\bigr),
 \qquad m_0=1,4.
\]
The empty configuration supplies zero.  Finally let
\(\ell\to\infty\) and \(\eta\downarrow0\).
This proves the pointwise limit.  Local uniform convergence follows
again from finite convex convergence on \(\Omega\).
\end{proof}

\subsection{The ideal transform and its maximizing densities}

Let \(H_t=H_{\psi_t}\) and \(P_1^t=\T_Aq_t\).
We will need more than convergence of these functions: we must also
locate all limits of their maximizing densities.
For a finite convex function \(F\) on an open convex set, write
\(\partial F(\lambda)\) for the vectors \(v\) satisfying
\(F(y)\ge F(\lambda)+v\cdot(y-\lambda)\) throughout that set.
These subgradients exist at every point, are locally bounded, and
pass to limits under local uniform convergence.  The last two
claims follow by applying the supporting inequality to points on
a fixed small ball about the argument.

\begin{lemma}\label{lem:initial-transform}
Uniformly on compact subsets of \(\Omega\),
\begin{equation}\label{eq:ideal-three-planes}
 P_1^t(\beta,h)\longrightarrow
 P_*(\beta,h)
 :=\max\left(0,\ ph+\frac{\beta Ap^2}{2},\
                     4ph-9p+8\beta Ap^2\right).
\end{equation}
Moreover, if \(t_j\to\infty\), \((\beta_j,h_j)\to(\beta,h)\in\Omega\),
and \(d_j\) maximizes the expression defining
\(P_1^{t_j}(\beta_j,h_j)\), then the \(d_j\) are bounded.
Every subsequential limit belongs to \(\{0,p,4p\}\) and maximizes
the corresponding limiting expression.
\end{lemma}

\begin{proof}
Applying the first inequality in \eqref{eq:superstable} to the
pure core gives \eqref{eq:coercivity} with
\(\kappa=t^2b^3/8\) and \(B_0=t^2/2\).  Consequently
\[
 H_t(\beta,d)
 \ge \beta t\left(\frac{b^3d^2}{8}-\frac d2\right)
                       +\frac{d\log d-d}{t}.
\]
For large \(t\), the quadratic in this bound dominates the added term
\(\beta Ad^2/2\).  On a fixed compact parameter set, the transformed
objective is negative for \(d\ge D\), with \(D\) independent of all
large \(t\).  Since \(d=0\) gives zero, all maximizers lie in \([0,D]\).

The partial conjugate of \(q_*\) is
\[
 H_*(d)=
 \begin{cases}
  0,&0\le d\le p,\\
  3(d-p),&p\le d\le4p,\\
  +\infty,&\text{otherwise}.
 \end{cases}
\]
Consider a subsequence with \(d_j\to d\).
For every \(z\in\R\), Fenchel's inequality gives
\[
 P_1^{t_j}(\beta_j,h_j)
 \le q_{t_j}(\beta_j,z)+(h_j-z)d_j
                                      +\frac{\beta_jAd_j^2}{2}.
\]
Take the upper limit and then the infimum over \(z\).
The result is bounded by
\[
 hd-H_*(d)+\frac{\beta Ad^2}{2}.
\]
In particular \(d\in[0,4p]\), since the transform is nonnegative.
This objective is strictly convex on each of the two linearity
intervals of \(H_*\); its maximum is therefore the right side of
\eqref{eq:ideal-three-planes}.

For the matching lower bound, choose a maximizing
\(d_0\in\{0,p,4p\}\) for this limiting objective.
Its global optimality says that, for every \(d'\),
\begin{equation}\label{eq:strict-density-support}
 H_*(d')\ge H_*(d_0)
 +(h+\beta Ad_0)(d'-d_0)
 +\frac{\beta A}{2}(d'-d_0)^2.
\end{equation}
Thus \(z=h+\beta Ad_0\) exposes the unique conjugate density \(d_0\)
for \(q_*\).  Set \(z_j=h_j+\beta_jAd_0\), and choose a subgradient
\((s'_j,d'_j)\in\partial q_{t_j}(\beta_j,z_j)\).
The physical pressure is nondecreasing in \(h\), so \(d'_j\ge0\);
these are admissible densities for the transform.
By Proposition~\ref{prop:free-pressure} and
\eqref{eq:strict-density-support}, every cluster here is \((0,d_0)\).
Fenchel equality at \(d'_j\) therefore gives
\begin{align*}
 P_1^{t_j}(\beta_j,h_j)
 &\ge q_{t_j}(\beta_j,z_j)
       +(h_j-z_j)d'_j+\frac{\beta_jA(d'_j)^2}{2}\\
 &\longrightarrow P_*(\beta,h).
\end{align*}
The upper and lower limits agree.  Equality in the upper argument
forces every limiting maximizing density to maximize the limiting
objective, and strict convexity on its two intervals confines such
densities to \(\{0,p,4p\}\).  The argument allows moving parameters;
it proves the assertions, including local uniform convergence.
\end{proof}

\subsection{Joint supports and the initial triple}

We next turn the maximizing densities into full supporting vectors.
We first record the support identities needed for this step.  For a finite convex function \(F\) on \(\Omega\), let
\[
 G_F(\lambda)=
 \left\{\lim_{j\to\infty}\nabla F(\lambda_j):
 \lambda_j\to\lambda,\quad F\text{ is differentiable at }\lambda_j\right\}.
\]
Local Lipschitz continuity, almost-everywhere differentiability, and the
supporting inequalities give the reachable-gradient formula
\cite[Theorem~9.61]{RockafellarWets1998}
\begin{equation}\label{eq:reachable-gradient-formula}
 \partial F(\lambda)=\conv G_F(\lambda),
\end{equation}
where \(G_F(\lambda)\) is nonempty and compact.  Lemma~\ref{lem:reachable-local} proves this identity and the supporting
facts in the planar setting needed here.

For a finite concave function of \(\beta>0\), write
\(\partial_\beta^+\) for its superdifferential.
Thus \(e\in-\partial_\beta^+H(\beta_0,d)\) means
\[
 H(\beta,d)\le H(\beta_0,d)-e(\beta-\beta_0)
 \qquad(\beta>0).
\]
The following identity links these canonical slopes to joint
pressure supports.

\begin{lemma}[Joint Fenchel supports]\label{lem:joint-support}
Let \(Q\) be finite and convex on \(\Omega\), and let
\[
 H(\beta,d)=\sup_{h\in\R}\{hd-Q(\beta,h)\}.
\]
Fix \(d\ge0\) for which \(H(\,\cdot\,,d)\) is finite and concave.
If \((\beta_0,z)\in\Omega\) satisfies
\[
 H(\beta_0,d)=zd-Q(\beta_0,z),
\]
then
\begin{equation}\label{eq:joint-support-identity}
 \{e\in\R:(e,d)\in\partial Q(\beta_0,z)\}
       =-\partial_\beta^+H(\beta_0,d).
\end{equation}
\end{lemma}

\begin{proof}
For \(e\) in the right side and every \((\beta,h)\in\Omega\),
the defining conjugate inequality gives
\begin{align*}
 Q(\beta,h)
 &\ge hd-H(\beta,d)\\
 &\ge Q(\beta_0,z)+e(\beta-\beta_0)+d(h-z).
\end{align*}
This is the joint supporting inequality.
Conversely, a support \((e,d)\) gives
\[
 hd-Q(\beta,h)\le
 zd-Q(\beta_0,z)-e(\beta-\beta_0)
\]
for every \(h\).  Take the supremum in \(h\) to obtain the
supergradient inequality for \(H(\,\cdot\,,d)\).
\end{proof}

The next observation records how a density maximizing a quadratic transform
also supports the original pressure at a shifted field. It transfers the full
supporting vector, including its first coordinate, which density convergence
alone would not control.

\begin{lemma}[Supports under a quadratic transform]\label{lem:transform-support}
Let $t$ satisfy \eqref{eq:t-choice}, let $Q=Q_\varphi$ for a potential
in \eqref{eq:physical-class}, and let $0<a\le A+1$. Put $H=H_Q$ and
$F=\T_aQ$. Suppose that $d\ge0$ attains the supremum defining
$F(\beta,h)$ at $(\beta,h)\in\Omega$, and put $z=h+\beta ad$.
Then, for every $d'\ge0$,
\begin{equation}\label{eq:quadratic-density-support}
 H(\beta,d')\ge H(\beta,d)+z(d'-d)
                         +\frac{\beta a}{2}(d'-d)^2.
\end{equation}
The finite concave function $H(\,\cdot\,,d)$ has a supergradient at $\beta$.
For every $e\in-\partial_\beta^+H(\beta,d)$,
\begin{equation}\label{eq:transform-support-transfer}
 (e,d)\in\partial Q(\beta,z),\qquad
 \left(e+\frac{ad^2}{2},d\right)\in\partial F(\beta,h).
\end{equation}
If $F$ is differentiable at $(\beta,h)$, the second vector equals
$\nabla F(\beta,h)$.
\end{lemma}

\begin{proof}
Compare the transformed objective at $d$ with its value at $d'$ and complete
the square to obtain \eqref{eq:quadratic-density-support}. Dropping the
nonnegative quadratic shows that $z$ supports $H(\beta,\cdot)$ at $d$,
so Fenchel equality holds. Lemma~\ref{lem:joint-support} gives the first
inclusion in \eqref{eq:transform-support-transfer}. For the second, the
fixed-density branch
\[
 (\beta',h')\longmapsto h'd-H(\beta',d)+\frac{\beta'a d^2}{2}
\]
lies below $F$ and agrees with it at $(\beta,h)$. The canonical
supergradient $-e$ gives this branch the displayed supporting vector,
which therefore supports $F$ as well. At a differentiability point there
is only one supporting vector.
\end{proof}

The three planes in \eqref{eq:ideal-three-planes} meet at
\(\lambda_*=(1,-3/4)\).  Their gradients are
\begin{equation}\label{eq:triangle-vertices}
 v_1=(0,0),\qquad
 v_2=(Ap^2/2,p)=(3p/4,p),\qquad
 v_3=(8Ap^2,4p)=(12p,4p).
\end{equation}
They are noncollinear, since \(\det(v_2,v_3)=-9p^2\ne0\).
Fix their barycenter
\begin{equation}\label{eq:fixed-center}
 g=\frac{v_1+v_2+v_3}{3}
   =\left(\frac{17p}{4},\frac{5p}{3}\right).
\end{equation}
Choose \(\tau>0\) small enough that the closed balls
\[
 U_i=\overline B(v_i,4\tau),\qquad i=1,2,3,
\]
have strictly ordered disjoint projections onto the density coordinate,
with the second projection contained in \((0,\infty)\), and have the
following properties.  Every triangle with one vertex in each \(U_i\)
is noncollinear and contains a common ball
\(\overline B(g,\gamma)\), for some \(\gamma>0\); also \(g\) lies in the
convex hull of no two of these balls.  Such a choice follows from the
strict interior position of \(g\) in the original triangle and compactness.
Set
\[
 r_0=\frac14,\qquad W=\overline B(\lambda_*,r_0)\subset\Omega .
\]
These constants will remain fixed throughout the construction.
Figure~\ref{fig:robust-triple} illustrates the common-ball property and
the separated density projections.

\begin{figure}[htbp]
\centering
\begin{tikzpicture}[x=1.45cm,y=1.05cm,>=Stealth,
                    every node/.style={font=\small}]
 \coordinate (Vone) at (0,0);
 \coordinate (Vtwo) at (1,2);
 \coordinate (Vthree) at (5.5,4);
 \coordinate (G) at (2.166667,2);
 \coordinate (Aone) at (.08,.07);
 \coordinate (Atwo) at (.90,2.07);
 \coordinate (Athree) at (5.55,3.87);
 \draw[->] (-.45,0) -- (5.95,0) node[right] {$s$};
 \draw[->] (0,-.45) -- (0,4.65) node[above] {$d$};
 \fill[black!3] (Vone) -- (Vtwo) -- (Vthree) -- cycle;
 \draw[black!55] (Vone) -- (Vtwo) -- (Vthree) -- cycle;
 \draw[black!65,dashed] (Aone) -- (Atwo) -- (Athree) -- cycle;
 \foreach \point in {Vone,Vtwo,Vthree}{
  \filldraw[fill=orange!15,draw=orange!65!black] (\point) circle[radius=.17];
 }
 \foreach \point in {Vone,Vtwo,Vthree}{
  \fill (\point) circle[radius=1pt];
 }
 \foreach \point in {Aone,Atwo,Athree}{
  \fill[black!65] (\point) circle[radius=1.2pt];
 }
 \node[below left] at (Vone) {$U_1$};
 \node[above left] at (Vtwo) {$U_2$};
 \node[above left] at (Vthree) {$U_3$};
 \filldraw[fill=blue!18,draw=blue!65!black] (G) circle[radius=.12];
 \fill (G) circle[radius=1.1pt];
 \node[above left] at (G) {$g$};
 \draw[blue!65!black] (2.30,1.94) -- (3.55,1.05);
 \node[blue!65!black,anchor=north,align=center] at (3.55,1.03)
   {$\overline B(g,\gamma)$\\same ball for every choice};
 \draw[black!25] (6.4,-.30) -- (6.4,4.35);
 \foreach \height in {0,2,4}{
  \draw[orange!65!black,line width=2.4pt]
    (6.4,{\height-.17}) -- (6.4,{\height+.17});
  \draw[black!25,densely dotted] (5.9,\height) -- (6.4,\height);
 }
 \node[anchor=west] at (6.52,0) {$\pi_2U_1$};
 \node[anchor=west] at (6.52,2) {$\pi_2U_2$};
 \node[anchor=west] at (6.52,4) {$\pi_2U_3$};
\end{tikzpicture}
\caption{Schematic geometry of the three slope neighborhoods. The two
triangles illustrate different choices of one point from each $U_i$.
The choice of the neighborhoods guarantees
$\overline B(g,\gamma)\subset\conv\{u_1,u_2,u_3\}$ for every
$u_i\in U_i$, while $g$ belongs to the convex hull of no two neighborhoods.
Their density projections $\pi_2U_i$, where $\pi_2(s,d)=d$, are disjoint. The actual supporting
groups lie in $d\ge0$, although $U_1$ extends below that line.}
\label{fig:robust-triple}
\end{figure}


\begin{proposition}[Initial supporting triple]\label{prop:initial-triple}
For all sufficiently large \(t\), there are
\(\lambda_1\in B(\lambda_*,r_0/16)\) and nonempty compact convex sets
\(K_i^1\subset\overline B(v_i,\tau/2)\), \(i=1,2,3\), such that
\begin{equation}\label{eq:initial-triple}
 g\in\partial P_1^t(\lambda_1)
       =\conv\bigl(K_1^1\cup K_2^1\cup K_3^1\bigr).
\end{equation}
All vectors in the \(K_i^1\) have nonnegative density coordinate.
In particular \(\partial P_1^t(\lambda_1)\) contains
\(\overline B(g,\gamma)\).
\end{proposition}

\begin{proof}
The function \(P_*(y)-g\cdot y\) has its unique minimum at
\(\lambda_*\).  In fact
\[
 P_*(y)-P_*(\lambda_*)-g\cdot(y-\lambda_*)
 =\max_i(v_i-g)\cdot(y-\lambda_*)\ge\gamma_0|y-\lambda_*|
\]
for some \(\gamma_0>0\), because the triangle of the \(v_i\)
contains a ball about \(g\).
Minimize \(P_1^t(y)-g\cdot y\) over
\(\overline B(\lambda_*,r_0/16)\).
Lemma~\ref{lem:initial-transform} shows that all minimizers lie in
the interior for large \(t\) and converge to \(\lambda_*\) as
\(t\to\infty\).  Choose one and call it \(\lambda_1^t\).
Interior minimality gives \(g\in\partial P_1^t(\lambda_1^t)\).
It remains to locate the slopes generating this subdifferential.

We claim that any limit of gradients
\(\nabla P_1^{t_j}(\beta_j,h_j)\), at differentiability points with
\(t_j\to\infty\) and \((\beta_j,h_j)\to\lambda_*\), belongs to
\(\{v_1,v_2,v_3\}\).  Choose a maximizing density \(d_j\) at each
such point. Apply Lemma~\ref{lem:transform-support} with $Q=q_{t_j}$
and $a=A$, choosing
$e_j\in-\partial_\beta^+H_{t_j}(\beta_j,d_j)$.
With $z_j=h_j+\beta_jAd_j$, it gives
\begin{equation}\label{eq:initial-gradient}
 \nabla P_1^{t_j}(\beta_j,h_j)
            =\left(e_j+\frac{Ad_j^2}{2},d_j\right),
\end{equation}
and
\begin{equation}\label{eq:initial-joint-support}
 (e_j,d_j)\in\partial q_{t_j}(\beta_j,z_j).
\end{equation}
The shifted support controls the original first-coordinate slope $e_j$
before the explicit increment $Ad_j^2/2$ is added.
The maximizing densities are bounded by
Lemma~\ref{lem:initial-transform}; hence the \(z_j\) stay in a
compact interval.  The locally uniform limit \(q_*\) is independent
of \(\beta\), so every cluster of the subgradients in
\eqref{eq:initial-joint-support} has first coordinate zero.
Therefore \(e_j\to0\).  The same lemma puts every density cluster
in \(\{0,p,4p\}\), and \eqref{eq:initial-gradient} proves the claim.

Subgradients of \(P_1^t\) near \(\lambda_*\) are bounded uniformly
for large \(t\), by \eqref{eq:ideal-three-planes}.
If \(G_{P_1^t}(\lambda_1^t)\) were not contained in
\(\bigcup_i\overline B(v_i,\tau/2)\) for arbitrarily large \(t\),
choose such a sequence of reachable gradients.  Approximate each
by a gradient at a differentiability point within distance \(1/t\)
of \(\lambda_1^t\), with gradient error tending to zero.
A bounded subsequence then contradicts the preceding claim.
Thus, for all sufficiently large \(t\), the three intersections
\[
 G_{P_1^t}(\lambda_1^t)\cap\overline B(v_i,\tau/2)
\]
cover the reachable-gradient set.  Let \(K_i^1\) be their convex
hulls.  They are compact and remain inside their respective balls.
Formula~\eqref{eq:reachable-gradient-formula} gives the equality
in \eqref{eq:initial-triple}.  None of the three sets is empty:
otherwise \(g\) would lie in the convex hull of two of the larger
balls \(U_i\).  Their density coordinates are nonnegative because
\(P_1^t\) is nondecreasing in \(h\).

Finally, choose one point in each \(K_i^1\).  Their triangle contains
\(\overline B(g,\gamma)\) by the fixed geometry of the \(U_i\), and
is contained in \(\partial P_1^t(\lambda_1^t)\).
\end{proof}

Fix from now on one \(t\) satisfying \eqref{eq:t-choice} and
Proposition~\ref{prop:initial-triple}, and suppress its superscript
when writing \(P_1=\T_Aq_t\) and \(\lambda_1\).
We have obtained a corner for an ideal transform.  A physical
finite-range approximation can smooth that corner, so uniform
approximation alone will not finish the proof.  The next section
shows how a further small mean-field attraction restores the
three separated groups of supporting slopes while moving their
common supporting point only a controlled distance.
